%%%%%%%%%%%%%%%%%%%%%%%%%%%%%%%%%%%%%%%%%%%%%%%%%%%%%%%%%%%%%%
%% Web GIS Laboratory 3: District Dashboard
%% Institute of Geographical Information Systems, NUST
%%
%% Layout follows "A Psychology Lab Manual" by John Paul Minda
%% (Overleaf, CC BY 4.0): article class, fullpage margins,
%% Helvetica throughout, parskip, blue links, no colour boxes.
%%
%% COMPILER: pdfLaTeX
%%%%%%%%%%%%%%%%%%%%%%%%%%%%%%%%%%%%%%%%%%%%%%%%%%%%%%%%%%%%%%

\documentclass{article}
\usepackage{fullpage}

\renewcommand{\familydefault}{\sfdefault}
\usepackage[scaled=1]{helvet}
\usepackage[helvet]{sfmath}
\everymath={\sf}

\usepackage{parskip}
\usepackage[colorlinks=true, allcolors=blue]{hyperref}

%% Required by this document but not by the original template.
%% Encoding first: without T1, < and > inside code print as
%% inverted punctuation and straight quotes turn curly.
\usepackage{iftex}
\ifPDFTeX
  \usepackage[T1]{fontenc}
  \usepackage[utf8]{inputenc}
\fi
\usepackage{textcomp}
\usepackage{upquote}
\usepackage{xcolor}
\usepackage{listings}
\usepackage{booktabs}
\usepackage{tabularx}
\usepackage{array}
\usepackage{enumitem}

\setlist[enumerate]{itemsep=.25em,topsep=.5em}
\setlist[itemize]{itemsep=.25em,topsep=.5em}
\renewcommand{\arraystretch}{1.25}
\newcolumntype{L}[1]{>{\raggedright\arraybackslash}p{#1}}

%% Inline code. \cd for short tokens, \cdb for long paths and URLs,
%% which must be allowed to break across lines.
\urlstyle{tt}
\newcommand{\cd}[1]{\texttt{#1}}
\newcommand{\cdb}[1]{\nolinkurl{#1}}

%% Code blocks: a plain rule, no shading, in keeping with the template
\lstdefinestyle{plain}{
  basicstyle=\ttfamily\small,
  frame=single,
  rulecolor=\color{black!35},
  framesep=7pt,
  xleftmargin=7pt,
  xrightmargin=7pt,
  breaklines=true,
  columns=fullflexible,
  keepspaces=true,
  showstringspaces=false,
  upquote=true,
  aboveskip=1em,
  belowskip=1em,
}
\lstset{style=plain}
\lstnewenvironment{code}{\lstset{style=plain}}{}

%% Caption for a code block. Deliberately not \paragraph, which is a
%% run-in heading and would be overlapped by the listing frame.
\newcommand{\codecap}[1]{%
  \par\vspace{.9em}\noindent\textbf{#1}\par\vspace{-.6em}}

%% A block that must not be skimmed past: ruled above and below,
%% with a bold label. Used for warnings and checkpoints.
\newenvironment{keypoint}[1]
  {\par\vspace{.7em}\noindent\rule{\linewidth}{.6pt}\par\vspace{.15em}
   \noindent\textbf{#1}\par\nobreak\vspace{-.3em}}
  {\par\vspace{.15em}\noindent\rule{\linewidth}{.6pt}\par\vspace{.7em}}

\title{Web GIS Laboratory 3\\[.3em]\large District Dashboard:
       layers, controls and interaction}
\author{Dr.\ Shahid Nawaz Khan\\
        Institute of Geographical Information Systems\\
        School of Civil and Environmental Engineering\\
        National University of Sciences and Technology, Islamabad}
\date{MS Remote Sensing and GIS}

\setcounter{tocdepth}{1}

\begin{document}
\maketitle

\begin{tabularx}{\textwidth}{@{}L{3.4cm} X@{}}
\toprule
Duration    & 3 hours \\
Covers      & Lecture: Leaflet, building the map in the browser \\
Builds on   & Lab 2, District Explorer \\
Software    & QGIS, Visual Studio Code, a browser, Git \\
Deliverable & A published dashboard and a short report \\
\bottomrule
\end{tabularx}

\tableofcontents


%==============================================================
\section*{Before you begin}
\addcontentsline{toc}{section}{Before you begin}

\subsection*{What you will build}

In Lab 2 you drew one layer and changed it with one drop-down. In this
lab you turn that map into a small dashboard, using most of what the
lecture covered:

\begin{itemize}
  \item a choice of three basemaps and a layer control to switch them;
  \item districts shaded by \textbf{population density}, with a panel
        that reports the district under the pointer;
  \item population drawn a second way, as \textbf{circles} sized by the
        number of people, in their own pane above the polygons;
  \item map events that can be switched on and off;
  \item an extension that measures what SVG and canvas drawing cost.
\end{itemize}

\subsection*{Learning outcomes}

By the end of this laboratory you will be able to:

\begin{enumerate}
  \item Configure a map through its options, and constrain its view.
  \item Offer several basemaps and overlays through a layer control, and
        meet the attribution each one requires.
  \item Explain why a choropleth must show a rate, and show a count
        with proportional symbols instead.
  \item Use all four hooks of \cd{L.geoJSON}: \cd{style},
        \cd{onEachFeature}, \cd{filter} and \cd{pointToLayer}.
  \item Build hover highlighting with \cd{resetStyle}, and a custom
        control that updates as the pointer moves.
  \item Control draw order with a pane.
  \item Register and remove event handlers, and explain why \cd{off}
        needs a named function.
  \item Measure the difference between SVG and canvas rendering.
\end{enumerate}

\subsection*{Structure and timing}

Parts 1 to 7 are the core. Parts 8 and 9 are extensions. If you run short
of time, complete Part 10 first so that your work is published.

\vspace{.6em}
\begin{tabularx}{\textwidth}{@{}c X r@{}}
\toprule
\textbf{Part} & \textbf{What you do} & \textbf{Time} \\
\midrule
1  & Set up the project                          & 10 min \\
2  & Make a point layer in QGIS                  & 15 min \\
3  & The page, the map and its basemaps          & 20 min \\
4  & A thematic layer, done properly             & 25 min \\
5  & Interaction: highlight and an info panel    & 25 min \\
6  & Proportional symbols in their own pane      & 25 min \\
7  & Map events, on and off                      & 15 min \\
\midrule
8  & Filter by province \textit{(extension)}     & 15 min \\
9  & SVG or canvas \textit{(extension)}          & 15 min \\
\midrule
10 & Publish and submit                          & 15 min \\
\bottomrule
\end{tabularx}

\subsection*{What you need}

\begin{tabularx}{\textwidth}{@{}L{4.2cm} X@{}}
\toprule
\textbf{Item} & \textbf{Requirement} \\
\midrule
Lab 2           & Finished and pushed. You reuse its data and its ideas \\
Data            & \cd{districts.geojson} from Lab 2, with \cd{pop\_2023}
                  and \cd{area\_km2} \\
QGIS            & Version 3.28 or later \\
Visual Studio Code & With Live Server \\
Browser         & Chrome or Edge, for the developer tools \\
\bottomrule
\end{tabularx}

\begin{keypoint}{This lab is longer on reasoning than on typing}
Most of the code is short. What matters is why each piece is written
the way it is: why density and not population, why a pane, why a named
function. The questions at the end ask about those decisions, not about
syntax.
\end{keypoint}


%==============================================================
\section{Set up the project}

Work in the same repository you used for Lab 2, so that Lab 2 stays
published. Create a new folder for this lab inside it:

\begin{code}
cd C:\webgis-lab
mkdir lab3
mkdir lab3\data
copy data\districts.geojson lab3\data\districts.geojson
code .
\end{code}

In \cd{lab3}, create \cd{index.html}, \cd{style.css} and \cd{app.js}.
When you publish, this lab will appear at your Lab 2 address followed by
\cd{/lab3/}.

\begin{keypoint}{Open the right folder in Live Server}
Live Server serves the folder VS Code has open. Open
\cd{lab3/index.html} and start Live Server from it, and check that the
address bar ends in \cd{/lab3/index.html}. The data paths in this lab
are relative to that file.
\end{keypoint}

\begin{keypoint}{Checkpoint 1}
A screenshot of VS Code showing \cd{lab3} with its three files and
\cd{data/districts.geojson}.
\end{keypoint}


%==============================================================
\section{Make a point layer in QGIS}

Population will be drawn twice: as a density on the polygons, and as a
circle per district. Circles need a point for each district.

\subsection{Point on surface, not centroid}

Load \cd{lab3/data/districts.geojson} into QGIS. Then run both of these
tools on it, from \textbf{Vector}, \textbf{Geometry Tools}:

\begin{tabularx}{\textwidth}{@{}L{3.6cm} X@{}}
\toprule
\textbf{Tool} & \textbf{Places the point} \\
\midrule
Centroids         & At the centre of mass. For a curved or hollow shape
                    this can fall outside the polygon \\
Point on Surface  & Guaranteed to fall inside the polygon \\
\bottomrule
\end{tabularx}

\vspace{.8em}
Find the district whose centroid lies outside its own boundary. There is
one in this dataset. Zoom to it and take a screenshot showing both
points.

\paragraph{Why it matters}
A symbol drawn outside its district is attributed by the reader to the
neighbouring one. On a dashboard nobody checks; they read it.

\subsection{Export}

Export the \textbf{Point on Surface} result as GeoJSON with these
settings, and record its size and feature count.

\begin{tabularx}{\textwidth}{@{}L{4.6cm} X@{}}
\toprule
\textbf{Setting} & \textbf{Value} \\
\midrule
File name            & \cd{lab3/data/district\_points.geojson} \\
CRS                  & EPSG:4326 \\
\cd{COORDINATE\_PRECISION} & 5 \\
\bottomrule
\end{tabularx}

\vspace{.8em}
The points inherit every attribute of their district, including
\cd{pop\_2023}. Confirm in the attribute table that there are 160
features.

\begin{keypoint}{Checkpoint 2}
The screenshot of the district whose centroid falls outside it, with
both points visible, and the size and feature count of
\cd{district\_points.geojson}.
\end{keypoint}


%==============================================================
\section{The page, the map and its basemaps}

\subsection{index.html}

\codecap{lab3/index.html, complete}
\begin{code}
<!DOCTYPE html>
<html lang="en">
<head>
  <meta charset="utf-8">
  <meta name="viewport" content="width=device-width, initial-scale=1">
  <title>District Dashboard</title>

  <link rel="stylesheet"
        href="https://unpkg.com/leaflet@1.9.4/dist/leaflet.css">
  <link rel="stylesheet" href="style.css">
</head>
<body>

  <header>
    <h1>District Dashboard</h1>
    <p class="subtitle">Population density and population of
      Pakistan's districts, 2023 census</p>
  </header>

  <div id="panel">
    <label for="province">Province</label>
    <select id="province">
      <option value="all">All provinces</option>
    </select>
    <span id="count">loading...</span>

    <label class="toggle">
      <input type="checkbox" id="coords">
      Show coordinates on click
    </label>
  </div>

  <div id="map"></div>

  <script src="https://unpkg.com/leaflet@1.9.4/dist/leaflet.js"></script>
  <script src="app.js"></script>
</body>
</html>
\end{code}

\subsection{style.css}

The first half repeats Lab 2. The second half styles the two controls
you will draw on the map.

\codecap{lab3/style.css, complete}
\begin{code}
body { font-family: system-ui, sans-serif; margin: 0;
       color: #23373b; background: #f4f6f6; }
header { padding: 18px 24px 10px; }
h1 { margin: 0; font-size: 24px; }
.subtitle { margin: 4px 0 0; color: #5b6a6e; }

#panel { padding: 10px 24px 14px; display: flex; flex-wrap: wrap;
         align-items: center; gap: 12px; }
#count { color: #5b6a6e; font-size: 14px; }
.toggle { margin-left: auto; font-size: 14px; }

#map { height: calc(100vh - 150px); min-height: 420px; width: 100%; }

.info, .legend { background: white; padding: 8px 10px;
                 border-radius: 4px; font-size: 13px; line-height: 1.45;
                 box-shadow: 0 1px 4px rgba(0, 0, 0, .25); }
.info { min-width: 200px; }
.info h4 { margin: 0 0 4px; font-size: 14px; }
.info .muted { color: #6b7a7e; }

.legend i { width: 16px; height: 16px; float: left;
            margin-right: 8px; opacity: .85; }
.legend .circles { margin-top: 8px; padding-top: 6px;
                   border-top: 1px solid #e3e6e7; }
.legend svg { display: block; margin-top: 4px; }
\end{code}

\subsection{The map, through its options}

In Lab 2 you created the map and then called \cd{setView}. Here every
setting goes into the options object instead.

\codecap{lab3/app.js, first block}
\begin{code}
const map = L.map('map', {
  center: [30.3753, 69.3451],
  zoom: 5,
  minZoom: 5,
  maxZoom: 14,
  maxBounds: [[21.5, 58.5], [39.0, 80.5]],
  maxBoundsViscosity: 1.0
});
\end{code}

\begin{tabularx}{\textwidth}{@{}L{4.2cm} X@{}}
\toprule
\textbf{Option} & \textbf{Why this value} \\
\midrule
\cd{minZoom: 5}  & The whole country fits; further out shows only
                   neighbours \\
\cd{maxZoom: 14} & District boundaries simplified to about 100 m stop
                   looking right well before street level \\
\cd{maxBounds}   & Pakistan with about two degrees to spare on each side \\
\bottomrule
\end{tabularx}

\subsection{Three basemaps and a layer control}

\codecap{lab3/app.js, continued}
\begin{code}
const osm = L.tileLayer('https://tile.openstreetmap.org/{z}/{x}/{y}.png', {
  maxZoom: 19,
  attribution: '&copy; OpenStreetMap contributors'
});

const positron = L.tileLayer(
  'https://{s}.basemaps.cartocdn.com/light_all/{z}/{x}/{y}{r}.png', {
  maxZoom: 19,
  subdomains: 'abcd',
  attribution: '&copy; OpenStreetMap contributors &copy; CARTO'
});

const imagery = L.tileLayer(
  'https://server.arcgisonline.com/ArcGIS/rest/services/World_Imagery/MapServer/tile/{z}/{y}/{x}', {
  maxZoom: 19,
  attribution: 'Tiles &copy; Esri'
});

positron.addTo(map);

const layerControl = L.control.layers(
  {
    'Light (Carto Positron)': positron,
    'OpenStreetMap': osm,
    'Satellite (Esri)': imagery
  },
  {},
  { collapsed: false }
).addTo(map);

L.control.scale({ imperial: false }).addTo(map);
\end{code}

The second argument to \cd{L.control.layers} is the overlays. It is
empty for now; Parts 4 and 6 add to it. Keep the reference in
\cd{layerControl}, because you need it to do that.

\paragraph{Read the templates closely}
Compare the three URL templates. Two of them end \cd{\{x\}/\{y\}}. Esri's
ends \cd{\{y\}/\{x\}}. Leaflet fills in whatever placeholders the
template contains, in whatever order they appear, so each provider's
template must be copied exactly. \cd{\{s\}} is replaced by one of the
\cd{subdomains} letters and \cd{\{r\}} by \cd{@2x} on high-density
screens.

\paragraph{Why Positron is the default}
A thematic layer is about to be drawn on top. A busy basemap competes
with it; a quiet grey one does not.

\begin{keypoint}{Attribution is a licence condition}
Each \cd{attribution} string is shown in the corner of the map, and it
changes when you switch basemap. Do not remove it or the attribution
control to tidy the map: the OpenStreetMap licence requires the credit.
\end{keypoint}

\begin{keypoint}{Checkpoint 3}
Open the Network tab, switch to each basemap in turn, and take one
screenshot showing an Esri tile request with its URL visible. In your
report, write down which number in that URL is the row and which is the
column.
\end{keypoint}


%==============================================================
\section{A thematic layer, done properly}

\subsection{Why density, not population}

In Lab 2 you shaded districts by population. That map has a flaw. A
large district can hold many people simply because it is large, so the
shading largely repeats the districts' areas. A choropleth colours
\textbf{areas}, so the value behind the colour must be one that does not
grow with area: a rate, a ratio or a density.

Population is still worth showing. It is a count, and counts are shown
with symbols whose size follows the count. That is Part 6.

\begin{keypoint}{The rule}
Choropleth: densities, rates, percentages. Proportional symbols:
counts. Never colour polygons by a count.
\end{keypoint}

\subsection{Style by density}

\codecap{lab3/app.js, continued}
\begin{code}
function show(value) {
  if (value === null || value === undefined || value === '') {
    return 'not recorded';
  }
  return typeof value === 'number' ? value.toLocaleString() : value;
}

function densityOf(p) {
  if (p.pop_2023 === null || p.pop_2023 === undefined) return null;
  return p.pop_2023 / p.area_km2;
}

function colourFor(density) {
  if (density === null) return '#c9cfd1';
  if (density > 1000) return '#08306b';
  if (density >  400) return '#2171b5';
  if (density >  150) return '#4292c6';
  if (density >   50) return '#9ecae1';
  return '#deebf7';
}

function styleFor(feature) {
  return {
    color: '#ffffff',
    weight: 1,
    fillColor: colourFor(densityOf(feature.properties)),
    fillOpacity: 0.8
  };
}
\end{code}

Density is not stored in the file. \cd{densityOf} works it out from two
attributes that are, every time a district is styled. With these breaks
the 135 districts with census data fall 20, 18, 34, 41 and 22 into the
five classes.

\subsection{A group the layer control can hold}

The layer control needs one object to switch on and off. But in Part 8
the district layer is rebuilt every time the province changes, and a
control pointing at a layer that has been thrown away switches nothing.
The answer is a group that stays put while its contents change.

\codecap{lab3/app.js, continued}
\begin{code}
const districtGroup = L.featureGroup().addTo(map);
layerControl.addOverlay(districtGroup, 'Population density');
\end{code}

\subsection{Load the data and draw}

\codecap{lab3/app.js, continued}
\begin{code}
let allDistricts = null;
let allPoints = null;
let districts = null;

function draw(province) {
  const keep = function (feature) {
    return province === 'all' || feature.properties.province === province;
  };

  districtGroup.clearLayers();

  districts = L.geoJSON(allDistricts, {
    filter: keep,
    style: styleFor
  });
  districtGroup.addLayer(districts);

  const n = districts.getLayers().length;
  if (n > 0) {
    map.fitBounds(districts.getBounds());
  }
  document.getElementById('count').textContent = n + ' districts';
}

async function loadJSON(path) {
  const response = await fetch(path);
  if (!response.ok) {
    throw new Error(path + ': status ' + response.status);
  }
  return response.json();
}

async function loadData() {
  try {
    [allDistricts, allPoints] = await Promise.all([
      loadJSON('data/districts.geojson'),
      loadJSON('data/district_points.geojson')
    ]);
  } catch (err) {
    console.error('Could not load the data.', err);
    document.getElementById('count').textContent = 'data failed to load';
    return;
  }
  draw('all');
}

loadData();
\end{code}

\paragraph{The \texttt{filter} hook}
In Lab 2 you filtered the array of features yourself and handed Leaflet
the result. Here \cd{L.geoJSON} does it: \cd{filter} is called once per
feature, and a feature for which it returns \cd{false} is never drawn.
\cd{keep} reads \cd{province} from the call to \cd{draw} that created
it, so the same function serves every province.

\paragraph{Two files, one wait}
\cd{Promise.all} starts both requests at once and waits for both. Two
\cd{await}s in a row would fetch the second file only after the first
had arrived.

\subsection{A legend for the density classes}

\codecap{lab3/app.js, continued}
\begin{code}
const legend = L.control({ position: 'bottomright' });

legend.onAdd = function () {
  const div = L.DomUtil.create('div', 'legend');
  const breaks = [0, 50, 150, 400, 1000];
  const labels = ['under 50', '50 to 150', '150 to 400',
                  '400 to 1,000', 'over 1,000'];

  div.innerHTML = '<strong>People per km<sup>2</sup></strong><br>';
  for (let i = 0; i < breaks.length; i++) {
    div.innerHTML +=
      '<i style="background:' + colourFor(breaks[i] + 1) + '"></i> ' +
      labels[i] + '<br>';
  }
  div.innerHTML +=
    '<i style="background:' + colourFor(null) + '"></i> no census data';

  L.DomEvent.disableClickPropagation(div);
  return div;
};

legend.addTo(map);
\end{code}

\begin{keypoint}{Checkpoint 4}
A screenshot of the density map with the legend, and with the layer
control open showing the three basemaps and the one overlay.
\end{keypoint}


%==============================================================
\section{Interaction: highlight and an info panel}

\subsection{A custom control that updates}

A popup costs the user a click and covers the map. An info panel in the
corner costs nothing: it reports whatever is under the pointer.

\codecap{lab3/app.js, continued}
\begin{code}
const info = L.control({ position: 'topright' });

info.onAdd = function () {
  this._div = L.DomUtil.create('div', 'info');
  L.DomEvent.disableClickPropagation(this._div);
  L.DomEvent.disableScrollPropagation(this._div);
  this.update();
  return this._div;
};

info.update = function (p) {
  const zoom = '<div class="muted">Zoom level ' + map.getZoom() + '</div>';
  if (!p) {
    this._div.innerHTML = '<h4>District</h4>' +
      '<div class="muted">Hover over a district</div>' + zoom;
    return;
  }
  const d = densityOf(p);
  this._div.innerHTML =
    '<h4>' + p.district + '</h4>' +
    p.province + '<br>' +
    'Population: ' + show(p.pop_2023) + '<br>' +
    'Area: ' + show(p.area_km2) + ' km<sup>2</sup><br>' +
    'Density: ' + (d === null ? 'not recorded'
       : Math.round(d).toLocaleString() + ' per km<sup>2</sup>') +
    zoom;
};

info.addTo(map);
\end{code}

\paragraph{\texttt{update} is yours}
Leaflet calls \cd{onAdd} once. \cd{update} is an ordinary method you
added to the control; nothing calls it until your own code does. Called
with no argument it shows the resting message.

\begin{keypoint}{Why the two \texttt{disable} calls}
The panel sits inside the map container. Without
\cd{disableClickPropagation}, clicking it also clicks the map beneath;
without \cd{disableScrollPropagation}, scrolling over it zooms the map.
You will see the first one matter in Part 7.
\end{keypoint}

\subsection{Highlight on hover}

\codecap{lab3/app.js, continued}
\begin{code}
function highlight(e) {
  e.target.setStyle({ weight: 3, color: '#1a202c', fillOpacity: 0.95 });
  e.target.bringToFront();
  info.update(e.target.feature.properties);
}

function unhighlight(e) {
  districts.resetStyle(e.target);
  info.update();
}

function zoomToDistrict(e) {
  map.fitBounds(e.target.getBounds(), { padding: [20, 20] });
}

function bindDistrict(feature, layer) {
  layer.bindTooltip(feature.properties.district, { sticky: true });
  layer.on({
    mouseover: highlight,
    mouseout: unhighlight,
    click: zoomToDistrict
  });
}
\end{code}

Now hand \cd{bindDistrict} to \cd{L.geoJSON} inside \cd{draw}:

\begin{code}
  districts = L.geoJSON(allDistricts, {
    filter: keep,
    style: styleFor,
    onEachFeature: bindDistrict
  });
\end{code}

\begin{tabularx}{\textwidth}{@{}L{3.8cm} X@{}}
\toprule
\textbf{Piece} & \textbf{Why} \\
\midrule
\cd{e.target}      & The district that fired the event. One shared
                     handler serves all 160 \\
\cd{bringToFront}  & Otherwise the thick outline is half hidden under
                     the neighbouring polygons \\
\cd{resetStyle}    & Re-runs \cd{styleFor} for that district. You never
                     store the original colour \\
\cd{districts}     & \cd{resetStyle} belongs to the \cd{L.geoJSON} layer,
                     which is why \cd{draw} keeps a reference to it \\
\bottomrule
\end{tabularx}

\begin{keypoint}{Checkpoint 5}
A screenshot with one district highlighted, its tooltip showing, and the
info panel reporting its population and density.
\end{keypoint}


%==============================================================
\section{Proportional symbols in their own pane}

\subsection{A pane above the polygons}

Polygons and circle markers are both vector paths, so by default both go
into \cd{overlayPane} and whichever is drawn later covers the other.
Hovering a district calls \cd{bringToFront}, which would lift it over
the circles. A pane with a higher z-index settles the order once.

\codecap{lab3/app.js, continued}
\begin{code}
map.createPane('symbols');
map.getPane('symbols').style.zIndex = 450;
\end{code}

\cd{overlayPane} is 400 and \cd{markerPane} 600, so 450 puts the circles
above every polygon and below markers, tooltips and popups.

\subsection{Size the circles}

\codecap{lab3/app.js, continued}
\begin{code}
function radiusFor(population) {
  return Math.sqrt(population) / 160;
}

function circleFor(feature, latlng) {
  return L.circleMarker(latlng, {
    pane: 'symbols',
    radius: radiusFor(feature.properties.pop_2023),
    color: '#7c2d12',
    weight: 1,
    fillColor: '#f97316',
    fillOpacity: 0.7
  }).bindTooltip(
    feature.properties.district + ': ' +
    feature.properties.pop_2023.toLocaleString()
  );
}

function hasPopulation(feature) {
  return feature.properties.pop_2023 !== null;
}
\end{code}

\paragraph{Why the square root}
The eye compares the \textbf{area} of a circle, and area grows with the
square of the radius. If the radius followed population, a district with
four times the people would look sixteen times larger. Taking the square
root first makes area follow population. Lahore, at about 13 million,
gets a radius of about 22 pixels.

\paragraph{\texttt{pointToLayer} and \texttt{filter} together}
Without \cd{pointToLayer}, every GeoJSON point becomes a default pin.
\cd{hasPopulation} drops the 25 districts with no census value before
anything is drawn, so \cd{circleFor} never meets a \cd{null}.

\subsection{Draw them, and add them to the control}

Below \cd{districtGroup}, add a second group:

\codecap{lab3/app.js, continued}
\begin{code}
const circleGroup = L.featureGroup().addTo(map);
layerControl.addOverlay(circleGroup, 'Population');
\end{code}

Then in \cd{draw}, clear it alongside the districts, and fill it
straight after \cd{districtGroup.addLayer(districts)}:

\begin{code}
  circleGroup.clearLayers();

  circleGroup.addLayer(L.geoJSON(allPoints, {
    filter: function (f) { return keep(f) && hasPopulation(f); },
    pointToLayer: circleFor
  }));
\end{code}

\subsection{Add the sizes to the legend}

Inside \cd{legend.onAdd}, just before the
\cd{disableClickPropagation} line, add three reference circles drawn
with the same \cd{radiusFor} as the map, so the legend cannot disagree
with it:

\begin{code}
  const sizes = [10000000, 4000000, 1000000];
  const rMax = radiusFor(sizes[0]);
  const pad = 6;
  const base = pad + 2 * rMax;
  let svg = '<svg width="' + (2 * rMax + 75) +
            '" height="' + (base + 2) + '">';
  for (const s of sizes) {
    const r = radiusFor(s);
    const top = base - 2 * r;
    svg += '<circle cx="' + (rMax + 1) + '" cy="' + (base - r) +
           '" r="' + r + '" fill="#f97316" fill-opacity=".2"' +
           ' stroke="#7c2d12"/>' +
           '<line x1="' + (rMax + 1) + '" y1="' + top +
           '" x2="' + (2 * rMax + 14) + '" y2="' + top +
           '" stroke="#7c2d12" stroke-width=".6"/>' +
           '<text x="' + (2 * rMax + 17) + '" y="' + (top + 3.5) +
           '" font-size="10">' + (s / 1000000) + ' million</text>';
  }
  svg += '</svg>';
  div.innerHTML +=
    '<div class="circles"><strong>Population</strong>' + svg + '</div>';
\end{code}

\subsection{Metres or pixels}

Add one more overlay, to see the difference between the two circle
constructors from the lecture:

\codecap{lab3/app.js, continued}
\begin{code}
const nust = [33.6423, 72.9906];
const nustGroup = L.featureGroup([
  L.circle(nust, { radius: 25000, color: '#0f766e', weight: 2,
                   fillOpacity: 0.08 }),
  L.circleMarker(nust, { radius: 6, color: '#0f766e',
                         fillColor: '#0f766e', fillOpacity: 1 })
]).bindTooltip('NUST H-12: a 25 km circle and a 6 pixel circleMarker');
layerControl.addOverlay(nustGroup, 'NUST, 25 km');
\end{code}

Turn it on in the layer control and zoom from level 5 to level 10. Note
what happens to each circle on screen, and measure the 25 km circle
against the scale bar.

\begin{keypoint}{Checkpoint 6}
Two screenshots: the full dashboard with circles above the density
shading and the legend showing both parts; and the NUST overlay at two
different zoom levels.
\end{keypoint}


%==============================================================
\section{Map events, on and off}

\subsection{A handler you can remove}

The checkbox in the panel switches a click handler on and off. While it
is on, clicking the map shows the coordinates of that point.

\codecap{lab3/app.js, continued}
\begin{code}
function showCoordinates(e) {
  L.popup()
    .setLatLng(e.latlng)
    .setContent(e.latlng.lat.toFixed(4) + ', ' + e.latlng.lng.toFixed(4))
    .openOn(map);
}

document.getElementById('coords')
  .addEventListener('change', function (event) {
    if (event.target.checked) {
      map.on('click', showCoordinates);
    } else {
      map.off('click', showCoordinates);
    }
  });
\end{code}

\begin{keypoint}{Why \texttt{showCoordinates} has a name}
\cd{map.off} removes a handler by finding the very function you passed
to \cd{map.on}. Suppose you had written this instead:

\begin{code}
map.on('click', function (e) { ... });
\end{code}

That function has no name and no variable holding it, so nothing you
could pass to \cd{off} would match it. The handler could never be
removed.
\end{keypoint}

\paragraph{Two listeners, two systems}
\cd{addEventListener} is the browser's own event system, used on page
elements. \cd{map.on} is Leaflet's, used on maps and layers. They look
alike and do the same job for different objects.

\subsection{Keep the panel's zoom level current}

\codecap{lab3/app.js, continued}
\begin{code}
map.on('zoomend', function () {
  info.update();
});
\end{code}

\cd{zoomend} fires once when a zoom finishes. \cd{zoom} would fire
continuously during the animation.

\subsection{Try these}

\begin{enumerate}
  \item Tick the checkbox and click the sea south of Karachi. Then click
        a district. Both handlers run: write down what you see.
  \item With the checkbox ticked, click the legend. No popup appears.
        In the legend's \cd{onAdd}, remove the line that calls
        \cd{disableClickPropagation}, reload, and click the legend again.
        Put the line back.
  \item Untick the checkbox and confirm clicks on the sea do nothing.
\end{enumerate}

\begin{keypoint}{Checkpoint 7}
A screenshot of a coordinate popup on the sea, and your notes from the
three steps above.
\end{keypoint}


%==============================================================
\section{Filter by province}
\textit{Extension exercise}

\cd{draw} already takes a province and filters by it. All that is
missing is the drop-down.

\codecap{lab3/app.js, continued}
\begin{code}
function fillProvinceList() {
  const select = document.getElementById('province');
  const provinces = [...new Set(
    allDistricts.features.map(f => f.properties.province)
  )].sort();
  for (const name of provinces) {
    const option = document.createElement('option');
    option.value = name;
    option.textContent = name;
    select.appendChild(option);
  }
}

document.getElementById('province')
  .addEventListener('change', function (event) {
    draw(event.target.value);
  });
\end{code}

Call \cd{fillProvinceList()} in \cd{loadData}, just before
\cd{draw('all')}.

\paragraph{What to check}
Select Gilgit Baltistan. The count reads 14 districts but no circles
appear. Explain why in your report, and say whether the dashboard should
tell the user so.

\begin{keypoint}{Checkpoint 8}
A screenshot with one province selected and the layer control open,
showing that the overlays still switch the rebuilt layers on and off.
\end{keypoint}


%==============================================================
\section{SVG or canvas}
\textit{Extension exercise}

\subsection{Make the renderer switchable}

Add this as the very first line of \cd{app.js}:

\begin{code}
const useCanvas = new URLSearchParams(window.location.search).has('canvas');
\end{code}

Add \cd{preferCanvas: useCanvas} to the options of \cd{L.map}. Opening
\cd{index.html?canvas} now draws every vector onto a canvas, and
\cd{index.html} on its own still uses SVG.

\subsection{Time the drawing}

Make the first line of \cd{draw}

\begin{code}
  const t0 = performance.now();
\end{code}

and the last

\begin{code}
  console.log('draw(' + province + '): ' + n + ' districts in ' +
              (performance.now() - t0).toFixed(1) + ' ms, ' +
              (useCanvas ? 'canvas' : 'SVG'));
\end{code}

\subsection{Measure}

For each renderer, reload five times and record the time printed for
\cd{draw(all)}. Then, in the Elements tab, find \cd{leaflet-overlay-pane}
and count what it contains.

\vspace{.6em}
\begin{tabularx}{\textwidth}{@{}X r r@{}}
\toprule
 & \textbf{SVG} & \textbf{Canvas} \\
\midrule
\cd{draw(all)}, median of five, ms & \hspace{2cm} & \hspace{2cm} \\
Elements inside \cd{leaflet-overlay-pane} & & \\
Hover highlight still works? & & \\
Can you inspect one district in Elements? & & \\
\bottomrule
\end{tabularx}

\paragraph{Reading the result}
With 160 districts and 135 circles, the difference is small and may
even favour SVG. The lecture's guidance is that canvas starts to pay at
a few thousand features. Say in your report what you would expect for
the tehsil layer (administrative level 3), which has several times as
many features, each with more vertices, and why.

\begin{keypoint}{Checkpoint 9}
The completed table, and screenshots of the Elements tab in both modes.
\end{keypoint}


%==============================================================
\section{Publish and submit}

\subsection{Commit and push}

\begin{code}
cd C:\webgis-lab
git add -A
git commit -m "Lab 3: district dashboard"
git push
\end{code}

Wait for GitHub Pages to rebuild, then open
\cdb{https://<your-username>.github.io/webgis-lab-<roll>/lab3/} on your
phone. Check that Lab 2 still works at its old address.

\subsection{Write the report}

Create \cd{lab3/report.md} containing:

\begin{enumerate}
  \item Your checkpoint screenshots.
  \item The size and feature count of \cd{district\_points.geojson}, and
        the name of the district whose centroid falls outside it.
  \item The row and column in your Esri tile URL.
  \item Your notes from the three steps in Part 7.
  \item The table from Part 9, if you completed it.
  \item Your answers to the questions below.
\end{enumerate}

\subsection{Questions}

\begin{enumerate}
  \item Lab 2 shaded districts by population. Name one large, thinly
        populated district and one small, crowded one, and explain how
        each would mislead on the Lab 2 map.
  \item \cd{radiusFor} takes a square root. What would the map
        suggest about Lahore against a district of 3 million if it did
        not?
  \item Why does the layer control hold \cd{districtGroup} rather than
        the \cd{L.geoJSON} layer itself?
  \item What would go wrong if the circles were drawn in
        \cd{overlayPane} instead of the \cd{symbols} pane? Describe what
        the user would see when hovering.
  \item Why can \cd{map.off} remove \cd{showCoordinates}, when it could
        not remove an anonymous function passed to \cd{map.on}?
  \item On the NUST overlay, which circle would you use to show a 25 km
        evacuation zone, and which to mark the campus on a national map?
        Why?
\end{enumerate}

\subsection{Marking}

\begin{tabularx}{\textwidth}{@{}X r@{}}
\toprule
\textbf{Item} & \textbf{Marks} \\
\midrule
Published dashboard with basemaps, density shading and legend & 20 \\
Hover highlight, info panel and click to zoom                  & 15 \\
Proportional symbols in their own pane, with legend           & 15 \\
Coordinate click handler that switches on and off             & 10 \\
Point layer, with the centroid comparison                     & 10 \\
Answers to the six questions                                  & 25 \\
Checkpoint screenshots and commit history                     & 5 \\
\midrule
Extensions in Parts 8 and 9                                   & +10 bonus \\
\bottomrule
\end{tabularx}

\vspace{1em}
Submit two links on LMS: your repository and your live \cd{lab3} page.


%==============================================================
\section*{Appendix A: Troubleshooting}
\addcontentsline{toc}{section}{Appendix A: Troubleshooting}

\begin{tabularx}{\textwidth}{@{}L{4.6cm} X@{}}
\toprule
\textbf{Symptom} & \textbf{Cause and fix} \\
\midrule
\cd{data failed to load}, 404 in the console &
  Live Server was started from the wrong folder, or the file is not in
  \cd{lab3/data}. Check the address bar ends in \cd{/lab3/index.html} \\
Satellite basemap shows grey or the wrong place &
  The \cd{\{y\}/\{x\}} order in the Esri template was changed. Copy it
  exactly \\
Highlight never goes away &
  \cd{resetStyle} was called on something other than the
  \cd{L.geoJSON} layer, or \cd{districts} was never assigned in
  \cd{draw} \\
\cd{Cannot read properties of null (reading 'resetStyle')} &
  The same: \cd{draw} must assign \cd{districts = L.geoJSON(...)} \\
Circles hidden under polygons &
  \cd{pane: 'symbols'} is missing from the circle options, or the pane
  was created after the circles were drawn \\
\cd{Cannot read properties of null (reading 'toLocaleString')} &
  A district with no population reached \cd{circleFor}. Add the
  \cd{hasPopulation} filter \\
Layer control switches nothing after changing province &
  The control was given the \cd{L.geoJSON} layer instead of the group \\
Clicking the legend opens a coordinate popup &
  \cd{disableClickPropagation} is missing from that control \\
Coordinate popups cannot be turned off &
  The handler passed to \cd{map.on} is an anonymous function. Name it \\
Info panel zoom level lags behind &
  The \cd{zoomend} handler is missing \\
\bottomrule
\end{tabularx}

%==============================================================
\section*{Appendix B: Reference}
\addcontentsline{toc}{section}{Appendix B: Reference}

\subsection*{The four hooks of \texttt{L.geoJSON}, as used here}

\begin{tabularx}{\textwidth}{@{}L{3.4cm} X@{}}
\toprule
\textbf{Hook} & \textbf{In this lab} \\
\midrule
\cd{filter}        & \cd{keep}: one province, and \cd{hasPopulation} for
                     the circles \\
\cd{style}         & \cd{styleFor}: colour by density \\
\cd{onEachFeature} & \cd{bindDistrict}: tooltip and three handlers \\
\cd{pointToLayer}  & \cd{circleFor}: a circle marker instead of a pin \\
\bottomrule
\end{tabularx}

\subsection*{Default panes}

\begin{tabularx}{\textwidth}{@{}L{4cm} r X@{}}
\toprule
\textbf{Pane} & \textbf{z-index} & \textbf{Holds} \\
\midrule
\cd{tilePane}    & 200 & Basemap tiles \\
\cd{overlayPane} & 400 & Vector paths: polygons, lines, circles \\
\cd{symbols}     & 450 & This lab's population circles \\
\cd{shadowPane}  & 500 & Marker shadows \\
\cd{markerPane}  & 600 & Marker icons \\
\cd{tooltipPane} & 650 & Tooltips \\
\cd{popupPane}   & 700 & Popups \\
\bottomrule
\end{tabularx}

\subsection*{Optional, if you finish everything}

\begin{enumerate}
  \item Add Leaflet.markercluster and draw the 160 points as clustered
        markers on a separate overlay.
  \item Show the density value as a permanent tooltip on each district
        above zoom level 8, using \cd{zoomend} to add and remove them.
  \item Save the map's centre and zoom in \cd{localStorage} on
        \cd{moveend}, and restore them when the page loads.
  \item Replace the fixed density breaks with quantiles computed from
        the data, and update the legend to match.
\end{enumerate}

\end{document}
